% Options for packages loaded elsewhere
\PassOptionsToPackage{unicode}{hyperref}
\PassOptionsToPackage{hyphens}{url}
\documentclass[
  french,
]{article}
\usepackage{xcolor}
\usepackage[margin=2cm]{geometry}
\usepackage{amsmath,amssymb}
\setcounter{secnumdepth}{-\maxdimen} % remove section numbering
\usepackage{iftex}
\ifPDFTeX
  \usepackage[T1]{fontenc}
  \usepackage[utf8]{inputenc}
  \usepackage{textcomp} % provide euro and other symbols
\else % if luatex or xetex
  \usepackage{unicode-math} % this also loads fontspec
  \defaultfontfeatures{Scale=MatchLowercase}
  \defaultfontfeatures[\rmfamily]{Ligatures=TeX,Scale=1}
\fi
\usepackage{lmodern}
\ifPDFTeX\else
  % xetex/luatex font selection
\fi
% Use upquote if available, for straight quotes in verbatim environments
\IfFileExists{upquote.sty}{\usepackage{upquote}}{}
\IfFileExists{microtype.sty}{% use microtype if available
  \usepackage[]{microtype}
  \UseMicrotypeSet[protrusion]{basicmath} % disable protrusion for tt fonts
}{}
\makeatletter
\@ifundefined{KOMAClassName}{% if non-KOMA class
  \IfFileExists{parskip.sty}{%
    \usepackage{parskip}
  }{% else
    \setlength{\parindent}{0pt}
    \setlength{\parskip}{6pt plus 2pt minus 1pt}}
}{% if KOMA class
  \KOMAoptions{parskip=half}}
\makeatother
\usepackage{longtable,booktabs,array}
\usepackage{caption}
\captionsetup[table]{skip=6pt}
\newcounter{none} % for unnumbered tables
\usepackage{calc} % for calculating minipage widths
% Correct order of tables after \paragraph or \subparagraph
\usepackage{etoolbox}
\makeatletter
\patchcmd\longtable{\par}{\if@noskipsec\mbox{}\fi\par}{}{}
\makeatother
\ifLuaTeX
\usepackage[bidi=basic,shorthands=off]{babel}
\else
\usepackage[bidi=default,shorthands=off]{babel}
\fi
\ifLuaTeX
  \usepackage{selnolig} % disable illegal ligatures
\fi
\setlength{\emergencystretch}{3em} % prevent overfull lines
\providecommand{\tightlist}{%
  \setlength{\itemsep}{0pt}\setlength{\parskip}{0pt}}
\usepackage{bookmark}
\IfFileExists{xurl.sty}{\usepackage{xurl}}{} % add URL line breaks if available
\urlstyle{same}
% fallback for those not using the hyperref driver hyperxmp:
\makeatletter
\@ifundefined{xmpquote}{\newcommand{\xmpquote}[1]{#1}}{}
\makeatother
\hypersetup{
  pdflang={fr},
  hidelinks,
  pdfcreator={LaTeX via pandoc}}

\author{}
\date{}

\begin{document}

\section{Manuel de Référence Théorique et Mathématique - Solveur COVOF
2D}\label{manuel-de-ruxe9fuxe9rence-thuxe9orique-et-mathuxe9matique---solveur-covof-2d}

\subsection{Modélisation Mathématique des Écoulements Incompressibles
Monophasiques, Diphasiques et
Convectifs}\label{moduxe9lisation-mathuxe9matique-des-uxe9coulements-incompressibles-monophasiques-diphasiques-et-convectifs}

\begin{center}\rule{0.5\linewidth}{0.5pt}\end{center}

\subsection{Sommaire}\label{sommaire}

\begin{enumerate}
\def\labelenumi{\arabic{enumi}.}
\tightlist
\item
  \hyperref[1-introduction-et-principes-fondamentaux]{Introduction et
  Principes Fondamentaux}
\item
  \hyperref[2-uxe9quations-duxe9tat-et-moduxe8le-continu-one-fluid]{Équations
  d'État et Modèle Continu ``One-Fluid''}

  \begin{itemize}
  \tightlist
  \item
    2.1. Équations Générales de Navier-Stokes Incompressibles
  \item
    2.2. Propriétés Physiques Equivalentes et Lois de Mélange
  \item
    2.3. Couplage Thermoconvectif de Boussinesq
  \item
    2.4. Modélisation des Obstacles et Parois Complexes par Pénalisation
    de Brinkman
  \end{itemize}
\item
  \hyperref[3-muxe9thode-du-lagrangien-augmentuxe9-augmented-lagrangian]{Méthode
  du Lagrangien Augmenté (Augmented Lagrangian)}

  \begin{itemize}
  \tightlist
  \item
    3.1. Problématique du Point-Selle et de la Contrainte
    d'Incompressibilité
  \item
    3.2. Principe Variationnel et Pénalisation de la Divergence
  \item
    3.3. Algorithme d'Uzawa et Mise à Jour de la Pression
  \item
    3.4. Analyse Spectrale et Choix Optimal du Paramètre d'Augmentation
  \item
    3.5. Analogies Physiques : Compressibilité Artificielle et Vitesse
    du Son Numérique
  \end{itemize}
\item
  \hyperref[4-discruxe9tisation-spatiale-sur-grille-cartuxe9sienne-duxe9caluxe9e-mac]{Discrétisation
  Spatiale sur Grille Cartésienne Décalée (MAC)}

  \begin{itemize}
  \tightlist
  \item
    4.1. Topologie de Harlow-Welch et Positionnement des Inconnues
  \item
    4.2. Équation de Quantité de Mouvement en \(u\) (Faces Verticales)
  \item
    4.3. Équation de Quantité de Mouvement en \(v\) (Faces Horizontales)
  \item
    4.4. Discrétisation des Contraintes Visqueuses et Moyenne Harmonique
    Nodale
  \item
    4.5. Discrétisation du Terme de Pénalisation
    \(\nabla (\nabla \cdot \mathbf{u})\)
  \item
    4.6. Discrétisation des Termes Convectifs (Schémas Centré, Amont et
    TVD)
  \end{itemize}
\item
  \hyperref[5-discruxe9tisation-de-luxe9quation-duxe9nergie-thermique]{Discrétisation
  de l'Équation d'Énergie Thermique}

  \begin{itemize}
  \tightlist
  \item
    5.1. Formulation Conservative Implicite
  \item
    5.2. Conditions aux Limites Thermiques (Dirichlet et Neumann
    Adiabatique)
  \item
    5.3. Définition et Calcul du Nombre de Nusselt Local et Moyen
  \end{itemize}
\item
  \hyperref[6-moduxe9lisation-des-uxe9coulements-diphasiques-et-tension-de-surface]{Modélisation
  des Écoulements Diphasiques et Tension de Surface}

  \begin{itemize}
  \tightlist
  \item
    6.1. Suivi d'Interface par Fraction Volumique (Volume of Fluid -
    VOF)
  \item
    6.2. Schémas TVD à Limiteurs de Flux (Superbee, Minmod, Van Leer)
  \item
    6.3. Modèle Continu de Force Capillaire (Continuous Surface Force -
    CSF)
  \item
    6.4. Calcul Régularisé de la Normale et de la Courbure
  \item
    6.5. Formulation ``Balanced-Force'' et Suppression des Courants
    Parasites
  \end{itemize}
\item
  \hyperref[7-algorithmes-de-ruxe9solution-alguxe9brique-et-solveurs-de-krylov]{Algorithmes
  de Résolution Algébrique et Solveurs de Krylov}

  \begin{itemize}
  \tightlist
  \item
    7.1. Structure Matricielle du Système Linéaire Couplé
  \item
    7.2. Algorithme BiCGSTAB(2) (Stabilized Bi-Conjugate Gradient Order
    2)
  \item
    7.3. Algorithme GMRES(m) (Restarted Generalized Minimal Residual)
  \item
    7.4. Algorithme du Gradient Conjugué Préconditionné (PCG)
  \item
    7.5. Techniques de Préconditionnement (Jacobi, SSOR, ILU(0),
    Bloc-Jacobi MPI)
  \end{itemize}
\item
  \hyperref[8-stabilituxe9-numuxe9rique-et-crituxe8res-cfl]{Stabilité
  Numérique et Critères CFL}

  \begin{itemize}
  \tightlist
  \item
    8.1. Restriction Convective CFL
  \item
    8.2. Restriction Visqueuse
  \item
    8.3. Restriction Capillaire de Brackbill
  \item
    8.4. Algorithme d'Adaptation Automatique du Pas de Temps
  \end{itemize}
\item
  \hyperref[9-nomenclature-et-symboles-mathuxe9matiques]{Nomenclature et
  Symboles Mathématiques}
\item
  \hyperref[10-ruxe9fuxe9rences-bibliographiques]{Références
  Bibliographiques}
\end{enumerate}

\begin{center}\rule{0.5\linewidth}{0.5pt}\end{center}

\subsection{1. Introduction et Principes
Fondamentaux}\label{introduction-et-principes-fondamentaux}

Le code \textbf{COVOF 2D} est un solveur de Mécanique des Fluides
Numérique (CFD) à haute performance conçu pour la simulation directe des
écoulements incompressibles, laminaires ou instationnaires,
monophasiques, thermoconvectifs et diphasiques à surface libre.

Son architecture théorique repose sur trois piliers fondamentaux hérités
des développements du laboratoire \textbf{TREFLE / I2M (UMR CNRS 5295)}
et documentés dans la thèse de doctorat de référence
(\emph{THSE\_MEL.pdf}) ainsi que dans les publications sur le code
\emph{Référence Académique} (\emph{article\_rev.pdf}, \emph{Cahier de
Validation et Vérification CFD 2D}) :

\begin{enumerate}
\def\labelenumi{\arabic{enumi}.}
\tightlist
\item
  \textbf{La discrétisation spatiale sur grille décalée (Marker-And-Cell
  - MAC)} : inventée par Harlow et Welch (1965), elle garantit un
  couplage vitesse-pression intrinsèquement stable sans nécessiter de
  stabilisation artificielle de type Rhie-Chow, et élimine tout mode
  parasite de pression en damier (\emph{checkerboard modes}).
\item
  \textbf{La méthode du Lagrangien Augmenté (Augmented Lagrangian)} :
  développée initialement par Fortin \& Glowinski (1983) et adaptée à la
  CFD par Caltagirone \& Vincent (2001), elle transforme le problème de
  point-selle indéfini de Navier-Stokes en un système elliptique défini
  positif strictement coercif, évitant l'inversion d'une équation de
  Poisson de pression globale mal conditionnée.
\item
  \textbf{Le modèle continu à un seul fluide (One-Fluid Model) couplé à
  la méthode CSF de Brackbill} : il permet de représenter des
  écoulements à interfaces multiples et à très forts contrastes de masse
  volumique (\(\rho_1 / \rho_2 \sim 1000\)) et de viscosité
  (\(\mu_1 / \mu_2 \sim 100\)) en traitant l'interface comme une
  singularité volumique régulière.
\end{enumerate}

\begin{center}\rule{0.5\linewidth}{0.5pt}\end{center}

\subsection{2. Équations d'État et Modèle Continu
``One-Fluid''}\label{uxe9quations-duxe9tat-et-moduxe8le-continu-one-fluid}

\subsubsection{2.1. Équations Générales de Navier-Stokes
Incompressibles}\label{uxe9quations-guxe9nuxe9rales-de-navier-stokes-incompressibles}

Soit un domaine spatial borné \(\Omega \subset \mathbb{R}^2\) de
frontière \(\partial \Omega\), et un intervalle de temps \([0, T]\).
Dans le cadre du modèle à un seul fluide, les équations de Navier-Stokes
incompressibles à propriétés thermo-physiques variables s'écrivent sous
forme conservative locale :

\[\rho \left( \frac{\partial \mathbf{u}}{\partial t} + \nabla \cdot (\mathbf{u} \otimes \mathbf{u}) \right) = - \nabla p + \nabla \cdot \boldsymbol{\tau} + \rho \mathbf{g} + \mathbf{F}_b + \mathbf{F}_\sigma - \frac{\mu}{K} \mathbf{u}\]

\[\nabla \cdot \mathbf{u} = 0\]

où : - \(\mathbf{u} = (u, v)^T\) est le champ eulérien de vitesse fluide
\([\text{m/s}]\). - \(p\) est la pression scalaire hydrodynamique
\([\text{Pa}]\). - \(\rho(\mathbf{x}, t)\) est la masse volumique locale
\([\text{kg/m}^3]\). - \(\mu(\mathbf{x}, t)\) est la viscosité dynamique
locale \([\text{Pa}\cdot\text{s}]\). - \(\mathbf{g} = (g_x, g_y)^T\) est
le vecteur accélération de la pesanteur \([\text{m/s}^2]\). -
\(\mathbf{F}_b\) est la force de flottabilité thermique de Boussinesq
\([\text{N/m}^3]\). - \(\mathbf{F}_\sigma\) est la force volumique de
tension de surface capillaire \([\text{N/m}^3]\). -
\(-\frac{\mu}{K} \mathbf{u}\) est le terme de force de traînée poreuse
de Brinkman servant à pénaliser les frontières immergées
\([\text{N/m}^3]\). - \(\boldsymbol{\tau}\) est le tenseur des
contraintes visqueuses de déformation de Newton pour un fluide
incompressible :

\[\boldsymbol{\tau} = 2 \mu \mathbf{D}(\mathbf{u}) = \mu \left( \nabla \mathbf{u} + (\nabla \mathbf{u})^T \right)\]

En coordonnées cartésiennes bidimensionnelles \((x, y)\), les
composantes scalaires de \(\boldsymbol{\tau}\) s'explicitent comme :
\[\tau_{xx} = 2 \mu \frac{\partial u}{\partial x}, \quad \tau_{yy} = 2 \mu \frac{\partial v}{\partial y}, \quad \tau_{xy} = \tau_{yx} = \mu \left( \frac{\partial u}{\partial y} + \frac{\partial v}{\partial x} \right)\]

La divergence du tenseur visqueux dans les équations scalaires de
quantité de mouvement prend alors la forme exacte :
\[(\nabla \cdot \boldsymbol{\tau})_x = \frac{\partial \tau_{xx}}{\partial x} + \frac{\partial \tau_{xy}}{\partial y} = \frac{\partial}{\partial x}\left( 2\mu \frac{\partial u}{\partial x} \right) + \frac{\partial}{\partial y}\left( \mu \left( \frac{\partial u}{\partial y} + \frac{\partial v}{\partial x} \right) \right)\]

\[(\nabla \cdot \boldsymbol{\tau})_y = \frac{\partial \tau_{yx}}{\partial x} + \frac{\partial \tau_{yy}}{\partial y} = \frac{\partial}{\partial x}\left( \mu \left( \frac{\partial u}{\partial y} + \frac{\partial v}{\partial x} \right) \right) + \frac{\partial}{\partial y}\left( 2\mu \frac{\partial v}{\partial y} \right)\]

\subsubsection{2.2. Propriétés Physiques Équivalentes et Lois de
Mélange}\label{propriuxe9tuxe9s-physiques-uxe9quivalentes-et-lois-de-muxe9lange}

Dans un écoulement diphasique (ex: eau/air ou bulle montante), la
présence d'une phase \(k \in \{1, 2\}\) est repérée par une fonction
indicatrice de présence ou fraction volumique
\(C(\mathbf{x}, t) \in [0, 1]\) (\(C=1\) dans la phase 1, \(C=0\) dans
la phase 2).

Pour assurer la cohérence thermodynamique et la stricte conservation de
la quantité de mouvement et des contraintes tangentielles à la traversée
d'une interface diffuse :

\paragraph{1. Masse volumique équivalente (Moyenne
Arithmétique)}\label{masse-volumique-uxe9quivalente-moyenne-arithmuxe9tique}

La masse étant une quantité conservative additive par unité de volume :
\[\rho(C) = C \rho_1 + (1 - C) \rho_2\]

\paragraph{2. Viscosité dynamique équivalente (Moyenne
Harmonique)}\label{viscosituxe9-dynamique-uxe9quivalente-moyenne-harmonique}

Selon les travaux théoriques de Caltagirone et Vincent (2001) et la
thèse de doctorat (\emph{THSE\_MEL.pdf}), la contrainte de cisaillement
\(\tau = \mu \frac{\partial u}{\partial y}\) étant continue à travers
l'interface, la moyenne volumique de la vitesse de déformation implique
une moyenne harmonique de la viscosité :
\[\frac{1}{\mu(C)} = \frac{C}{\mu_1} + \frac{1 - C}{\mu_2} \quad \Longleftrightarrow \quad \mu(C) = \frac{\mu_1 \mu_2}{C \mu_2 + (1 - C) \mu_1}\]

Cette moyenne harmonique prévient l'apparition d'oscillations parasites
de contraintes et garantit que la phase la moins visqueuse lubrifie
adéquatement l'écoulement interfacial.

\subsubsection{2.3. Couplage Thermoconvectif de
Boussinesq}\label{couplage-thermoconvectif-de-boussinesq}

Lorsque les variations de température \(T(\mathbf{x}, t)\) sont
modérées, l'approximation de Boussinesq postule que les variations de
masse volumique sont négligeables dans tous les termes de l'équation de
Navier-Stokes, excepté dans le terme de gravité :
\[\rho(T) = \rho_0 \left( 1 - \beta (T - T_0) \right)\] où
\(\beta = -\frac{1}{\rho_0} \left( \frac{\partial \rho}{\partial T} \right)_p\)
est le coefficient de dilatation volumique thermique
\([\text{K}^{-1}]\), \(T_0\) la température de référence, et
\(\rho_0 = \rho(T_0)\).

Le terme de force de flottabilité s'écrit alors vectoriellement :
\[\mathbf{F}_b = -\rho_0 \beta (T - T_0) \mathbf{g}\]

Dans le cas canonique de la convection naturelle dans une cavité
différentiellement chauffée soumise à la gravité verticale
\(\mathbf{g} = (0, -g)^T\) :
\[F_{b, x} = 0, \quad F_{b, y} = \rho_0 \beta (T - T_0) g\]

\subsubsection{2.4. Modélisation des Obstacles et Parois Complexes par
Pénalisation de
Brinkman}\label{moduxe9lisation-des-obstacles-et-parois-complexes-par-puxe9nalisation-de-brinkman}

Pour représenter des géométries arbitraires (cylindres circulaires,
parois obliques, marches descendantes, cavités en forme de L) sur un
maillage cartésien régulier sans remaillage non structuré coûteux, COVOF
2D utilise la méthode de pénalisation de Brinkman (Angot et al.~1999,
Caltagirone 1994).

Le domaine solide \(\Omega_s \subset \Omega\) est assimilé à un milieu
poreux rigide de perméabilité intrinsèque quasi-nulle \(K \to 0\) :
\[\mathbf{F}_{pen} = - \frac{\mu}{K(\mathbf{x})} \mathbf{u}\] avec :
\[K(\mathbf{x}) = \begin{cases} +\infty & \text{si } \mathbf{x} \in \Omega_f \text{ (fluide)} \\ K_{solid} \ll 1 \quad (10^{-40}\text{ m}^2) & \text{si } \mathbf{x} \in \Omega_s \text{ (solide)} \end{cases}\]

Dans le fluide, \(\frac{\mu}{K} = 0\), et les équations de Navier-Stokes
standards sont rigoureusement retrouvées. Dans le solide, le terme
\(\frac{\mu}{K_{solid}} \mathbf{u}\) domine asymptotiquement tous les
autres termes, imposant numériquement :
\[\frac{\mu}{K_{solid}} \mathbf{u} = \mathbf{0} \quad \Longrightarrow \quad \mathbf{u}_{solid} \approx \mathbf{0}\]
avec une précision de l'ordre de \(O(K_{solid})\).

\begin{center}\rule{0.5\linewidth}{0.5pt}\end{center}

\subsection{3. Méthode du Lagrangien Augmenté (Augmented
Lagrangian)}\label{muxe9thode-du-lagrangien-augmentuxe9-augmented-lagrangian}

\subsubsection{3.1. Problématique du Point-Selle et de la Contrainte
d'Incompressibilité}\label{probluxe9matique-du-point-selle-et-de-la-contrainte-dincompressibilituxe9}

Dans un fluide incompressible, la pression n'est pas une variable
thermodynamique d'état déterminée par une loi des gaz, mais un
multiplicateur de Lagrange associé à la contrainte cinématique
\(\nabla \cdot \mathbf{u} = 0\).

Sous forme matricielle discrétisée standard, le système couplé
vitesse-pression s'écrit :
\[\begin{pmatrix} \mathbf{A} & \mathbf{B}^T \\ \mathbf{B} & \mathbf{0} \end{pmatrix} \begin{pmatrix} \mathbf{U} \\ \mathbf{P} \end{pmatrix} = \begin{pmatrix} \mathbf{F} \\ \mathbf{0} \end{pmatrix}\]
où \(\mathbf{A}\) représente l'opérateur de convection-diffusion
instationnaire, \(\mathbf{B}^T\) est la matrice du gradient de pression
(\(-\nabla\)), et \(\mathbf{B}\) la matrice de divergence
(\(\nabla \cdot\)).

Ce système présente un \textbf{bloc diagonal nul} \(\mathbf{0}\) en
pression. Il est indéfini et ne peut être résolu directement par des
méthodes de Krylov simples (Gradient Conjugué). De plus, les méthodes de
projection classiques (type Chorin-Temam) découplent vitesse et pression
en résolvant une équation de Poisson
\(\nabla^2 p = \frac{\rho}{\Delta t} \nabla \cdot \mathbf{u}^*\), ce qui
introduit des conditions aux limites artificielles sur la pression
(\(\frac{\partial p}{\partial n} = 0\)) responsables d'erreurs
numériques de couche limite d'ordre 1 en temps.

\subsubsection{3.2. Principe Variationnel et Pénalisation de la
Divergence}\label{principe-variationnel-et-puxe9nalisation-de-la-divergence}

La méthode du \textbf{Lagrangien Augmenté} résout cette difficulté en
ajoutant à l'équation de quantité de mouvement le gradient de la
divergence multiplié par un paramètre d'augmentation scalaire \(r > 0\)
:
\[\rho \frac{\mathbf{u}^{n+1} - \mathbf{u}^n}{\Delta t} + \rho (\mathbf{u}^n \cdot \nabla)\mathbf{u}^{n+1} - \nabla \cdot \left( 2 \mu \mathbf{D}(\mathbf{u}^{n+1}) \right) - r \nabla (\nabla \cdot \mathbf{u}^{n+1}) + \nabla p^{n+1} = \rho \mathbf{g} + \mathbf{F}_{ext}\]

Comme le fluide réel satisfait \(\nabla \cdot \mathbf{u} = 0\), le terme
\(-r \nabla (\nabla \cdot \mathbf{u}^{n+1})\) est \textbf{formellement
nul à la solution exacte}. Il n'introduit aucune altération physique de
la solution convergée !

\subsubsection{3.3. Algorithme d'Uzawa et Mise à Jour de la
Pression}\label{algorithme-duzawa-et-mise-uxe0-jour-de-la-pression}

En approchant \(p^{n+1}\) au début du pas de temps par \(p^n\), on
découple temporairement la vitesse et la pression via un algorithme de
point fixe d'Uzawa :

\paragraph{Étape 1 : Résolution de la Vitesse
Augmentée}\label{uxe9tape-1-ruxe9solution-de-la-vitesse-augmentuxe9e}

On résout pour \(\mathbf{u}^{n+1}\) le système elliptique :
\[\left( \frac{\rho}{\Delta t} \mathbf{I} + \rho (\mathbf{u}^n \cdot \nabla) - \nabla \cdot (2\mu \mathbf{D}) - r \nabla \text{div} \right) \mathbf{u}^{n+1} = \frac{\rho}{\Delta t} \mathbf{u}^n - \nabla p^n + \rho \mathbf{g} + \mathbf{F}_{ext}\]

\paragraph{Étape 2 : Actualisation d'Uzawa de la
Pression}\label{uxe9tape-2-actualisation-duzawa-de-la-pression}

La pression est immédiatement actualisée sans inversion de matrice :
\[p^{n+1} = p^n - r \nabla \cdot \mathbf{u}^{n+1}\]

Si l'incompressibilité n'est pas suffisamment satisfaite après un seul
pas (\(\|\nabla \cdot \mathbf{u}^{n+1}\|_\infty > \epsilon_{div}\)), on
réitère la boucle d'Uzawa :
\[p^{n+1, k+1} = p^{n+1, k} - r \nabla \cdot \mathbf{u}^{n+1, k+1}\]
Dans COVOF 2D, grâce au choix optimal de \(r\), une à deux itérations
d'Uzawa suffisent pour obtenir une divergence résiduelle
\(\|\nabla \cdot \mathbf{u}\| < 10^{-4}\text{ s}^{-1}\).

\subsubsection{3.4. Analyse Spectrale et Choix Optimal du Paramètre
d'Augmentation}\label{analyse-spectrale-et-choix-optimal-du-paramuxe8tre-daugmentation}

Soit l'opérateur \(\mathbf{L}_r = -\mu \nabla^2 - r \nabla \text{div}\).
Sur un domaine périodique, décomposons le champ de vitesse en une partie
solénoïdale \(\mathbf{u}_s\) (\(\nabla \cdot \mathbf{u}_s = 0\)) et une
partie irrotationnelle \(\mathbf{u}_d\)
(\(\nabla \times \mathbf{u}_d = \mathbf{0}\)) par le théorème de
décomposition de Helmholtz-Hodge :

\[\mathbf{L}_r \mathbf{u}_s = -\mu \nabla^2 \mathbf{u}_s\]
\[\mathbf{L}_r \mathbf{u}_d = -(\mu + r) \nabla^2 \mathbf{u}_d\]

Le paramètre \(r\) n'affecte donc \textbf{que la partie divergente} du
champ, qu'il atténue fortement. Caltagirone et Vincent (2001) ont
démontré que pour minimiser le conditionnement du système discret sur
une grille de pas d'espace \(h = \min(\Delta x, \Delta y)\), la valeur
optimale du paramètre \(r\) s'exprime sous la forme :

\[r_{opt} \approx \frac{\rho}{\Delta t} \left( \frac{L}{2\pi} \right)^2 \quad \text{ou empiriquement} \quad r_{opt} \sim 10^2 \text{ à } 10^4 \times \frac{\mu}{\rho}\]

Dans COVOF 2D, la valeur par défaut \(r = 100\text{ à }1000\) fournit
une convergence ultrarapide du solveur BiCGSTAB(2).

\subsubsection{3.5. Analogies Physiques : Compressibilité Artificielle
et Vitesse du Son
Numérique}\label{analogies-physiques-compressibilituxe9-artificielle-et-vitesse-du-son-numuxe9rique}

En prenant la divergence de l'équation de mise à jour de la pression :
\[\nabla \cdot \mathbf{u}^{n+1} = - \frac{1}{r} (p^{n+1} - p^n) \approx - \frac{\Delta t}{r} \frac{\partial p}{\partial t}\]
En comparant cette équation à l'équation de continuité d'un fluide
légèrement compressible :
\[\nabla \cdot \mathbf{u} = - \frac{1}{\rho c^2} \frac{\partial p}{\partial t}\]
on identifie immédiatement que la méthode du Lagrangien Augmenté
introduit une \textbf{vitesse du son numérique artificielle} :
\[c_{num} = \sqrt{\frac{r}{\rho \Delta t}}\]

Pour de grandes valeurs de \(r\), \(c_{num} \to \infty\), ce qui
correspond rigoureusement à la limite incompressible où les ondes
acoustiques se propagent instantanément.

\begin{center}\rule{0.5\linewidth}{0.5pt}\end{center}

\subsection{4. Discrétisation Spatiale sur Grille Cartésienne Décalée
(MAC)}\label{discruxe9tisation-spatiale-sur-grille-cartuxe9sienne-duxe9caluxe9e-mac}

\subsubsection{4.1. Topologie de Harlow-Welch et Positionnement des
Inconnues}\label{topologie-de-harlow-welch-et-positionnement-des-inconnues}

Le domaine de calcul \(\Omega = [0, L_x] \times [0, L_y]\) est
partitionné en \(N_x \times N_y\) cellules de contrôle rectangulaires de
dimensions \(\Delta x = L_x / N_x\) et \(\Delta y = L_y / N_y\).

\begin{verbatim}
         (i, j+1/2) : v[i, j+1]
              ▲
              │
  ┌───────────┴───────────┐ (i+1/2, j+1/2) : node (mu_node)
  │                       │
  │      p[i, j]          │
  │      C[i, j]          │
◄─┼──────T[i, j]──────────┼─► (i+1/2, j) : u[i+1, j]
  │      rho[i, j]        │
  │      mu[i, j]         │
  └───────────┬───────────┘
              │
              ▼
         (i, j-1/2) : v[i, j]
\end{verbatim}

\begin{itemize}
\tightlist
\item
  \textbf{Inconnues scalaires au centre de la cellule \((i, j)\)} :
  Pression \(p_{i,j}\), fraction volumique \(C_{i,j}\), température
  \(T_{i,j}\), masse volumique \(\rho_{i,j}\), viscosité dynamique
  moyenne \(\mu_{i,j}\).
\item
  \textbf{Vitesse normale \(u_{i, j}\) aux faces verticales} : située au
  point \((x_{i-1/2}, y_j)\), entre la maille \((i-1, j)\) et la maille
  \((i, j)\).
\item
  \textbf{Vitesse normale \(v_{i, j}\) aux faces horizontales} : située
  au point \((x_i, y_{j-1/2})\), entre la maille \((i, j-1)\) et la
  maille \((i, j)\).
\item
  \textbf{Viscosité nodale \(\mu_{node, i, j}\) aux sommets des mailles}
  : située au point \((x_{i-1/2}, y_{j-1/2})\).
\end{itemize}

\subsubsection{\texorpdfstring{4.2. Équation de Quantité de Mouvement en
\(u\) (Faces
Verticales)}{4.2. Équation de Quantité de Mouvement en u (Faces Verticales)}}\label{uxe9quation-de-quantituxe9-de-mouvement-en-u-faces-verticales}

Le volume de contrôle décalé associé à la composante \(u_{i,j}\) a pour
centre \((x_{i-1/2}, y_j)\) et pour volume
\(\Delta V_u = \Delta x \Delta y\).

L'intégration de l'équation scalaire en \(u\) sur \(\Delta V_u\) s'écrit
:
\[\iint_{\Delta V_u} \left[ \frac{\rho_u}{\Delta t} u^{n+1} - \frac{\partial \tau_{xx}}{\partial x} - \frac{\partial \tau_{xy}}{\partial y} - r \frac{\partial (\nabla \cdot \mathbf{u})}{\partial x} + \frac{\partial p}{\partial x} \right] dx dy = \iint_{\Delta V_u} \left[ \frac{\rho_u}{\Delta t} u^n + \rho_u g_x + F_{\sigma, x} - \rho_u (\mathbf{u} \cdot \nabla)u \right] dx dy\]

où \(\rho_u = \frac{\rho_{i-1, j} + \rho_{i, j}}{2}\).

\paragraph{\texorpdfstring{1. Terme normal visqueux
\(\frac{\partial \tau_{xx}}{\partial x}\)
:}{1. Terme normal visqueux \textbackslash frac\{\textbackslash partial \textbackslash tau\_\{xx\}\}\{\textbackslash partial x\} :}}\label{terme-normal-visqueux-fracpartial-tau_xxpartial-x}

\[\int_{y_{j-1/2}}^{y_{j+1/2}} \left( \tau_{xx}(x_i, y) - \tau_{xx}(x_{i-1}, y) \right) dy \approx \left[ 2 \mu_{i, j} \frac{u_{i+1, j} - u_{i, j}}{\Delta x} - 2 \mu_{i-1, j} \frac{u_{i, j} - u_{i-1, j}}{\Delta x} \right] \Delta y\]

\paragraph{\texorpdfstring{2. Terme tangentiel visqueux
\(\frac{\partial \tau_{xy}}{\partial y}\)
:}{2. Terme tangentiel visqueux \textbackslash frac\{\textbackslash partial \textbackslash tau\_\{xy\}\}\{\textbackslash partial y\} :}}\label{terme-tangentiel-visqueux-fracpartial-tau_xypartial-y}

\[\int_{x_{i-1}}^{x_i} \left( \tau_{xy}(x, y_{j+1/2}) - \tau_{xy}(x, y_{j-1/2}) \right) dx \approx \left[ \tau_{xy}(i, j+1/2) - \tau_{xy}(i, j-1/2) \right] \Delta x\]
avec :
\[\tau_{xy}(i, j+1/2) = \mu_{node}(i, j+1) \left( \frac{u_{i, j+1} - u_{i, j}}{\Delta y} + \frac{v_{i, j+1} - v_{i-1, j+1}}{\Delta x} \right)\]
\[\tau_{xy}(i, j-1/2) = \mu_{node}(i, j) \left( \frac{u_{i, j} - u_{i, j-1}}{\Delta y} + \frac{v_{i, j} - v_{i-1, j}}{\Delta x} \right)\]

\paragraph{\texorpdfstring{3. Terme de pénalisation du Lagrangien
Augmenté \(-r \frac{\partial (\nabla \cdot \mathbf{u})}{\partial x}\)
:}{3. Terme de pénalisation du Lagrangien Augmenté -r \textbackslash frac\{\textbackslash partial (\textbackslash nabla \textbackslash cdot \textbackslash mathbf\{u\})\}\{\textbackslash partial x\} :}}\label{terme-de-puxe9nalisation-du-lagrangien-augmentuxe9--r-fracpartial-nabla-cdot-mathbfupartial-x}

\[\int_{y_{j-1/2}}^{y_{j+1/2}} -r \left( (\nabla \cdot \mathbf{u})_{i, j} - (\nabla \cdot \mathbf{u})_{i-1, j} \right) dy = - r \left[ \left( \frac{u_{i+1, j} - u_{i, j}}{\Delta x} + \frac{v_{i, j+1} - v_{i, j}}{\Delta y} \right) - \left( \frac{u_{i, j} - u_{i-1, j}}{\Delta x} + \frac{v_{i-1, j+1} - v_{i-1, j}}{\Delta y} \right) \right] \Delta y\]

En regroupant les contributions en \(u\), le terme diagonal vaut :
\[\text{Diag}_{AL} = 2 r \frac{\Delta y}{\Delta x}\] et les voisins
Ouest et Est reçoivent chacun \(-r \frac{\Delta y}{\Delta x}\).

\paragraph{\texorpdfstring{4. Gradient de pression discret
\(-\frac{\partial p}{\partial x}\)
:}{4. Gradient de pression discret -\textbackslash frac\{\textbackslash partial p\}\{\textbackslash partial x\} :}}\label{gradient-de-pression-discret--fracpartial-ppartial-x}

\[-\int_{y_{j-1/2}}^{y_{j+1/2}} \left( p_{i, j} - p_{i-1, j} \right) dy = - (p_{i, j} - p_{i-1, j}) \Delta y\]

\subsubsection{\texorpdfstring{4.3. Équation de Quantité de Mouvement en
\(v\) (Faces
Horizontales)}{4.3. Équation de Quantité de Mouvement en v (Faces Horizontales)}}\label{uxe9quation-de-quantituxe9-de-mouvement-en-v-faces-horizontales}

Par symétrie, sur le volume de contrôle décalé centré en
\((x_i, y_{j-1/2})\) :
\[\rho_v = \frac{\rho_{i, j-1} + \rho_{i, j}}{2}\]

\paragraph{\texorpdfstring{1. Terme normal visqueux
\(\frac{\partial \tau_{yy}}{\partial y}\)
:}{1. Terme normal visqueux \textbackslash frac\{\textbackslash partial \textbackslash tau\_\{yy\}\}\{\textbackslash partial y\} :}}\label{terme-normal-visqueux-fracpartial-tau_yypartial-y}

\[\left[ 2 \mu_{i, j} \frac{v_{i, j+1} - v_{i, j}}{\Delta y} - 2 \mu_{i, j-1} \frac{v_{i, j} - v_{i, j-1}}{\Delta y} \right] \Delta x\]

\paragraph{\texorpdfstring{2. Terme tangentiel visqueux
\(\frac{\partial \tau_{yx}}{\partial x}\)
:}{2. Terme tangentiel visqueux \textbackslash frac\{\textbackslash partial \textbackslash tau\_\{yx\}\}\{\textbackslash partial x\} :}}\label{terme-tangentiel-visqueux-fracpartial-tau_yxpartial-x}

\[\left[ \mu_{node}(i+1, j) \left( \frac{v_{i+1, j} - v_{i, j}}{\Delta x} + \frac{u_{i+1, j} - u_{i+1, j-1}}{\Delta y} \right) - \mu_{node}(i, j) \left( \frac{v_{i, j} - v_{i-1, j}}{\Delta x} + \frac{u_{i, j} - u_{i, j-1}}{\Delta y} \right) \right] \Delta y\]

\paragraph{\texorpdfstring{3. Terme de pénalisation du Lagrangien
Augmenté \(-r \frac{\partial (\nabla \cdot \mathbf{u})}{\partial y}\)
:}{3. Terme de pénalisation du Lagrangien Augmenté -r \textbackslash frac\{\textbackslash partial (\textbackslash nabla \textbackslash cdot \textbackslash mathbf\{u\})\}\{\textbackslash partial y\} :}}\label{terme-de-puxe9nalisation-du-lagrangien-augmentuxe9--r-fracpartial-nabla-cdot-mathbfupartial-y}

\[- r \left[ \left( \frac{u_{i+1, j} - u_{i, j}}{\Delta x} + \frac{v_{i, j+1} - v_{i, j}}{\Delta y} \right) - \left( \frac{u_{i+1, j-1} - u_{i, j-1}}{\Delta x} + \frac{v_{i, j} - v_{i, j-1}}{\Delta y} \right) \right] \Delta x\]

\subsubsection{4.4. Discrétisation des Contraintes Visqueuses et Moyenne
Harmonique
Nodale}\label{discruxe9tisation-des-contraintes-visqueuses-et-moyenne-harmonique-nodale}

Au sommet \((i-1/2, j-1/2)\), la viscosité \(\mu_{node}(i, j)\) est
calculée par la moyenne harmonique des 4 mailles contiguës :
\[\mu_{node}(i, j) = \frac{4}{\frac{1}{\mu_{i-1, j-1}} + \frac{1}{\mu_{i, j-1}} + \frac{1}{\mu_{i-1, j}} + \frac{1}{\mu_{i, j}}}\]

Cette formule fondamentale permet de reproduire exactement les
conditions de saut interfacial de Navier-Stokes sans lissage artificiel.

\subsubsection{4.5. Discrétisation des Termes
Convectifs}\label{discruxe9tisation-des-termes-convectifs}

Le terme convectif \((\mathbf{u} \cdot \nabla) \phi\) pour une quantité
transportée \(\phi \in \{u, v\}\) est évalué selon trois schémas
commutables :

\paragraph{1. Schéma Centré d'Ordre 2
(Centered)}\label{schuxe9ma-centruxe9-dordre-2-centered}

\[u \left.\frac{\partial \phi}{\partial x}\right|_i = u_i \frac{\phi_{i+1} - \phi_{i-1}}{2\Delta x}\]
Non dissipatif, d'ordre 2 en espace, mais instable pour les nombres de
Reynolds de maille élevés
(\(Pe_{cell} = \frac{\rho |u| \Delta x}{\mu} > 2\)).

\paragraph{2. Schéma Amont d'Ordre 1
(Upwind)}\label{schuxe9ma-amont-dordre-1-upwind}

\[u \left.\frac{\partial \phi}{\partial x}\right|_i = \max(u_i, 0) \frac{\phi_i - \phi_{i-1}}{\Delta x} + \min(u_i, 0) \frac{\phi_{i+1} - \phi_i}{\Delta x}\]
Inconditionnellement stable, strictement monotone, mais introduit une
viscosité numérique diffusive d'ordre \(\frac{|u| \Delta x}{2}\).

\begin{center}\rule{0.5\linewidth}{0.5pt}\end{center}

\subsection{5. Discrétisation de l'Équation d'Énergie
Thermique}\label{discruxe9tisation-de-luxe9quation-duxe9nergie-thermique}

\subsubsection{5.1. Formulation Conservative
Implicite}\label{formulation-conservative-implicite}

L'équation gouvernant le champ de température \(T(\mathbf{x}, t)\) pour
un fluide de capacité thermique massique \(c_p\) et de conductivité
thermique \(k\) s'écrit :
\[\rho c_p \left( \frac{\partial T}{\partial t} + \nabla \cdot (\mathbf{u} T) \right) = \nabla \cdot (k \nabla T) + S_T\]

Sous forme discrète semi-implicite au centre de la maille \((i, j)\) :
\[\rho_{i,j} c_p \frac{T_{i,j}^{n+1} - T_{i,j}^n}{\Delta t} \Delta x \Delta y + \mathcal{C}_T(T^{n+1}) = \left[ k_{i+1/2, j} \frac{T_{i+1, j}^{n+1} - T_{i, j}^{n+1}}{\Delta x} - k_{i-1/2, j} \frac{T_{i, j}^{n+1} - T_{i-1, j}^{n+1}}{\Delta x} \right] \Delta y + \left[ k_{i, j+1/2} \frac{T_{i, j+1}^{n+1} - T_{i, j}^{n+1}}{\Delta y} - k_{i, j-1/2} \frac{T_{i, j}^{n+1} - T_{i, j-1}^{n+1}}{\Delta y} \right] \Delta x\]

Les conductivités aux faces sont calculées par moyenne harmonique :
\[k_{i+1/2, j} = \frac{2 k_{i,j} k_{i+1,j}}{k_{i,j} + k_{i+1,j}}\]

\subsubsection{5.2. Conditions aux Limites
Thermiques}\label{conditions-aux-limites-thermiques}

\begin{enumerate}
\def\labelenumi{\arabic{enumi}.}
\tightlist
\item
  \textbf{Dirichlet (Paroi isotherme \(T = T_w\))} : En utilisant une
  maille fictive fantôme (\emph{ghost cell}) d'indice \(i=0\) pour une
  paroi située en \(x=0\) :
  \[\frac{T_{ghost} + T_{1, j}}{2} = T_w \quad \Longrightarrow \quad T_{0, j} = 2 T_w - T_{1, j}\]
\item
  \textbf{Neumann (Paroi adiabatique
  \(\frac{\partial T}{\partial n} = 0\))} : \[T_{0, j} = T_{1, j}\]
\end{enumerate}

\subsubsection{5.3. Définition et Calcul du Nombre de Nusselt Local et
Moyen}\label{duxe9finition-et-calcul-du-nombre-de-nusselt-local-et-moyen}

Le transfert de chaleur adimensionnel est caractérisé par le
\textbf{Nombre de Nusselt} \(Nu\). À la paroi chaude verticale
(\(x = 0\)) maintenue à \(T_h\), en vis-à-vis d'une paroi froide
(\(x = L_x\)) à \(T_c\) : - \textbf{Nombre de Nusselt local} \(Nu(y)\) :
\[Nu(y) = - \left. \frac{\partial T^*}{\partial x^*} \right|_{x=0} = - \frac{L_x}{T_h - T_c} \left( \frac{T_{1, j} - T_w}{\Delta x / 2} \right)\]
- \textbf{Nombre de Nusselt moyen} \(\overline{Nu}\) :
\[\overline{Nu} = \frac{1}{L_y} \int_0^{L_y} Nu(y) dy = \frac{1}{N_y} \sum_{j=1}^{N_y} Nu(j)\]

\begin{center}\rule{0.5\linewidth}{0.5pt}\end{center}

\subsection{6. Modélisation des Écoulements Diphasiques et Tension de
Surface}\label{moduxe9lisation-des-uxe9coulements-diphasiques-et-tension-de-surface}

\subsubsection{6.1. Suivi d'Interface par Fraction Volumique (Volume of
Fluid -
VOF)}\label{suivi-dinterface-par-fraction-volumique-volume-of-fluid---vof}

La fraction volumique \(C(\mathbf{x}, t) \in [0, 1]\) vérifie l'équation
d'advection pure conservative :
\[\frac{\partial C}{\partial t} + \nabla \cdot (C \mathbf{u}) = 0\]

\subsubsection{6.2. Schémas TVD à Limiteurs de Flux (Superbee, Minmod,
Van
Leer)}\label{schuxe9mas-tvd-uxe0-limiteurs-de-flux-superbee-minmod-van-leer}

Pour empêcher la diffusion de l'interface tout en évitant toute
oscillation dispersive non physique (\(C < 0\) ou \(C > 1\)), COVOF 2D
met en œuvre la famille des schémas Total Variation Diminishing (TVD) :

\[F_{i+1/2} = u_{i+1/2} \left[ C_{up} + \frac{1}{2} (1 - |\nu|) \Phi(\theta) (C_{down} - C_{up}) \right]\]

où \(\nu = u_{i+1/2} \frac{\Delta t}{\Delta x}\) est le nombre de
Courant local, et
\(\theta = \frac{C_{up} - C_{up-1}}{C_{down} - C_{up}}\) est le rapport
des gradients locaux.

Les limiteurs disponibles sont : - \textbf{Superbee (Roe)} : le plus
compressif, optimal pour les interfaces diphasiques nettes :
\[\Phi_{Superbee}(\theta) = \max \left( 0, \max(\min(2\theta, 1), \min(\theta, 2)) \right)\]
- \textbf{Van Leer} : limiteur régulier et différentiable :
\[\Phi_{VanLeer}(\theta) = \frac{\theta + |\theta|}{1 + |\theta|}\] -
\textbf{Minmod} : le plus dissipatif et le plus robuste :
\[\Phi_{Minmod}(\theta) = \max(0, \min(1, \theta))\]

\subsubsection{6.3. Modèle Continu de Force Capillaire (Continuous
Surface Force -
CSF)}\label{moduxe8le-continu-de-force-capillaire-continuous-surface-force---csf}

Selon le modèle CSF de Brackbill et al.~(1992), l'effet de la tension
superficielle est converti en une force volumique concentrée dans la
zone de transition interfaciale :
\[\mathbf{F}_\sigma = \sigma \kappa \nabla C\] où \(\sigma\) est le
coefficient de tension superficielle \([\text{N/m}]\), et \(\kappa\) est
la courbure moyenne locale \([\text{m}^{-1}]\).

\subsubsection{6.4. Calcul Régularisé de la Normale et de la
Courbure}\label{calcul-ruxe9gularisuxe9-de-la-normale-et-de-la-courbure}

Le vecteur normal unitaire à l'interface est défini par :
\[\mathbf{n} = \frac{\nabla C}{\|\nabla C\|}\]

Pour éviter une division par zéro dans le fluide pur, une régularisation
est appliquée :
\[\mathbf{n}_{i, j} = \begin{cases} \frac{\nabla C_{i, j}}{\sqrt{\|\nabla C_{i, j}\|^2 + \epsilon^2}} & \text{si } \|\nabla C_{i, j}\| > \frac{10^{-3}}{\Delta x} \\ \mathbf{0} & \text{sinon} \end{cases}\]

La courbure est ensuite obtenue par la divergence conservative de la
normale aux centres des cellules :
\[\kappa_{i, j} = - \nabla \cdot \mathbf{n} = - \left( \frac{n_{x, i+1/2, j} - n_{x, i-1/2, j}}{\Delta x} + \frac{n_{y, i, j+1/2} - n_{y, i, j-1/2}}{\Delta y} \right)\]

\subsubsection{6.5. Formulation ``Balanced-Force'' et Suppression des
Courants
Parasites}\label{formulation-balanced-force-et-suppression-des-courants-parasites}

Sur une goutte statique au repos, l'équilibre de Navier-Stokes se réduit
rigoureusement à :
\[\nabla p = \mathbf{F}_\sigma = \sigma \kappa \nabla C\]

Si le gradient discret de pression \(\nabla p\) et le gradient discret
de fraction volumique \(\nabla C\) ne sont pas évalués
\textbf{exactement aux mêmes positions stencil} (aux faces décalées), un
déséquilibre numérique se crée, générant des tourbillons artificiels
appelés \textbf{courants parasites}.

Dans COVOF 2D, l'approximation ``Balanced-Force'' de François et
al.~(2006) est strictement respectée :
\[F_{\sigma, x, i, j} = \sigma \left( \frac{\kappa_{i-1, j} + \kappa_{i, j}}{2} \right) \left( \frac{C_{i, j} - C_{i-1, j}}{\Delta x} \right)\]
\[F_{\sigma, y, i, j} = \sigma \left( \frac{\kappa_{i, j-1} + \kappa_{i, j}}{2} \right) \left( \frac{C_{i, j} - C_{i, j-1}}{\Delta y} \right)\]

Cette formulation garantit que le saut de pression théorique de Laplace
: \[\Delta P = \frac{\sigma}{R}\] est exactement compensé par la force
capillaire discrète, maintenant les vitesses parasites résiduelles sous
le seuil de \(10^{-3}\text{ m/s}\).

\subsubsection{6.6. Limiteurs TVD et Diagramme NVD Avancés (Superbee,
Van Leer, Minmod, MUSCL,
HRIC)}\label{limiteurs-tvd-et-diagramme-nvd-avancuxe9s-superbee-van-leer-minmod-muscl-hric}

Le transport convectif de la fraction volumique \(C\) est modélisé par
l'équation conservative :
\[\frac{\partial C}{\partial t} + \nabla \cdot (C \mathbf{u}) = 0\]

Pour préserver la monotonie stricte (\(0 \le C \le 1\)) et la netteté de
l'interface sans diffusion excessive, COVOF 2D dispose de limiteurs de
flux formulés soit dans le formalisme TVD (Total Variation Diminishing,
Harten 1983) en fonction du ratio de pente \(\theta\), soit dans le
formalisme NVD (Normalized Variable Diagram, Leonard 1991).

Le flux convectif à travers une face \(i+1/2\) s'exprime par :
\[F_{i+1/2} = u_{i+1/2} \left[ C_{up} + \frac{1}{2} (1 - |\nu_{i+1/2}|) \Phi(\theta) (C_{down} - C_{up}) \right]\]
où \(\nu_{i+1/2} = u_{i+1/2} \frac{\Delta t}{\Delta x}\) est le nombre
de Courant local de face, \(C_{up}\) la valeur amont (upwind),
\(C_{down}\) la valeur aval (downwind), et \(\theta\) le rapport des
gradients amont/aval :
\[\theta = \frac{C_{up} - C_{up-1}}{C_{down} - C_{up}}\]

Dans le formalisme NVD de Leonard (1991), la variable normalisée s'écrit
: \[\tilde{C} = \frac{C - C_U}{C_A - C_U} = \frac{\theta}{\theta + 1}\]
où \(U\) désigne le nœud amont, \(A\) le nœud accepteur et \(D\) le nœud
donneur.

Les limiteurs implémentés dans COVOF 2D sont :

\begin{enumerate}
\def\labelenumi{\arabic{enumi}.}
\item
  \textbf{Superbee (Roe, 1985)} : Limiteur compressif d'extrême
  précision pour les interfaces diphasiques :
  \[\Phi_{Superbee}(\theta) = \max \left( 0, \max(\min(2\theta, 1), \min(\theta, 2)) \right)\]
\item
  \textbf{Van Leer (1979)} : Limiteur lisse et continu :
  \[\Phi_{VanLeer}(\theta) = \frac{\theta + |\theta|}{1 + |\theta|}\]
\item
  \textbf{Minmod (Roe, 1986)} : Limiteur le plus dissipatif et le plus
  stable : \[\Phi_{Minmod}(\theta) = \max(0, \min(1, \theta))\]
\item
  \textbf{MUSCL (Monotonic Upstream-Centered Scheme for Conservation
  Laws, Van Leer 1979)} :
  \[\Phi_{MUSCL}(\theta) = \max \left( 0, \min \left( 2\theta, \min \left( \frac{1 + \theta}{2}, 2 \right) \right) \right)\]
\item
  \textbf{HRIC (High-Resolution Interface Capturing, Muzaferija \& Peric
  1999)} : Formulé nativement dans le diagramme NVD avec correction
  locale du nombre de Courant \(Co\) :
  \[\tilde{C}_f^* = \begin{cases} 2 \tilde{C} & \text{si } 0 \le \tilde{C} < 0.5 \\ 1 & \text{si } 0.5 \le \tilde{C} \le 1 \\ \tilde{C} & \text{sinon} \end{cases}\]
  \[\tilde{C}_f = \sqrt{Co} \, \tilde{C}_f^* + (1 - \sqrt{Co}) \, \tilde{C} \quad (0 \le Co \le 1)\]
  Le coefficient de flux équivalent TVD est déduit par la transformation
  canonique :
  \[\Phi_{HRIC}(\theta) = \frac{2}{1 - Co} \left( \frac{\tilde{C}_f - \tilde{C}}{1 - \tilde{C}} \right) \quad (\tilde{C} \neq 1)\]
\end{enumerate}

\subsubsection{6.7. Méthodes Géométriques de Reconstruction d'Interface
(SLIC, PLIC, LVIRA,
ELVIRA)}\label{muxe9thodes-guxe9omuxe9triques-de-reconstruction-dinterface-slic-plic-lvira-elvira}

Dans une cellule diphasique mixte
(\(10^{-6} < C_{i, j} < 1 - 10^{-6}\)), l'interface réelle est approchée
localement par un segment de droite :
\[\mathbf{n} \cdot (\mathbf{x} - \mathbf{x}_0) = 0 \quad \iff \quad n_x x + n_y y + \alpha = 0\]
avec le vecteur normal unitaire sortant normalisé
\(\|\mathbf{n}\|_2 = \sqrt{n_x^2 + n_y^2} = 1\). Quatre algorithmes de
reconstruction sont implémentés dans COVOF 2D :

\paragraph{6.7.1. SLIC (Simple Line Interface Calculation, Noh \&
Woodward
1976)}\label{slic-simple-line-interface-calculation-noh-woodward-1976}

Approximation d'ordre 1 où le segment est aligné strictement avec l'un
des axes cartésiens. En analysant les gradients horizontaux et verticaux
voisins, la normale est fixée à :
\[\mathbf{n} \in \{ (1, 0), (-1, 0), (0, 1), (0, -1) \}\] Le paramètre
de position \(\alpha\) est ensuite déterminé de façon univoque par la
fraction volumique locale \(C_{i, j}\).

\paragraph{6.7.2. PLIC (Piecewise Linear Interface Calculation, Youngs
1982, Scardovelli \& Zaleski
1999)}\label{plic-piecewise-linear-interface-calculation-youngs-1982-scardovelli-zaleski-1999}

Approximation linéaire continue d'ordre 2. Le vecteur normal est estimé
par un gradient centré pondéré sur le stencil \(3 \times 3\) :
\[n_x = - \frac{1}{\Delta x} \left[ (C_{i+1, j+1} + 2 C_{i+1, j} + C_{i+1, j-1}) - (C_{i-1, j+1} + 2 C_{i-1, j} + C_{i-1, j-1}) \right]\]
\[n_y = - \frac{1}{\Delta y} \left[ (C_{i+1, j+1} + 2 C_{i, j+1} + C_{i-1, j+1}) - (C_{i+1, j-1} + 2 C_{i, j-1} + C_{i-1, j-1}) \right]\]
\[\mathbf{n} = \frac{\mathbf{n}}{\|\mathbf{n}\|_2}\]

\paragraph{6.7.3. LVIRA (Least-Squares Volume-of-Fluid Interface
Reconstruction Algorithm, Puckett
1991)}\label{lvira-least-squares-volume-of-fluid-interface-reconstruction-algorithm-puckett-1991}

Optimisation continue de l'angle d'interface \(\theta \in [0, 2\pi]\)
minimisant au sens des moindres carrés la différence entre les fractions
volumiques prescrites et les volumes tronqués sur les 9 cellules du
stencil :
\[\theta^* = \arg\min_{\theta} E(\theta) = \arg\min_{\theta} \sum_{k=-1}^1 \sum_{l=-1}^1 \left[ C_{i+k, j+l}^{num}(\mathbf{n}(\theta), \alpha) - C_{i+k, j+l} \right]^2\]
Dans COVOF 2D, cette minimisation non linéaire unimodale par morceaux
est résolue par la méthode de la section dorée (Golden Section Search).

\paragraph{6.7.4. ELVIRA (Efficient Least-Squares Volume-of-Fluid
Interface Reconstruction Algorithm, Pilliod \& Puckett
2004)}\label{elvira-efficient-least-squares-volume-of-fluid-interface-reconstruction-algorithm-pilliod-puckett-2004}

Algorithme haute performance qui évalue un ensemble discret de 6 pentes
candidates issues de fonctions de hauteur (Height Functions) colonnes et
lignes : - Pentes verticales :
\(m_{v1} = \frac{h_y(i+1) - h_y(i-1)}{2\Delta x}\),
\(m_{v2} = \frac{h_y(i+1) - h_y(i)}{\Delta x}\),
\(m_{v3} = \frac{h_y(i) - h_y(i-1)}{\Delta x}\) - Pentes horizontales :
\(m_{h1} = \frac{h_x(j+1) - h_x(j-1)}{2\Delta y}\),
\(m_{h2} = \frac{h_x(j+1) - h_x(j)}{\Delta y}\),
\(m_{h3} = \frac{h_x(j) - h_x(j-1)}{\Delta y}\) - Candidat de Youngs :
\(\mathbf{n}_{Youngs}\)

La pente candidate sélectionnée est celle qui minimise rigoureusement
l'erreur quadratique \(\sum |C_{num} - C_{exact}|^2\) sur le stencil
\(3 \times 3\). Cette méthode garantit un ordre de convergence
géométrique de 2 pour des interfaces régulières arbitrairement
inclinées.

\subsubsection{6.8. Découpe Polygonale de Sutherland-Hodgman et
Inversion de
Volume}\label{duxe9coupe-polygonale-de-sutherland-hodgman-et-inversion-de-volume}

Pour calculer la constante \(\alpha\) correspondant exactement à une
fraction volumique \(C_{i, j}\), COVOF 2D utilise une formulation
analytique exacte basée sur l'algorithme de Sutherland-Hodgman (1974) :
1. La cellule cartésienne
\([x_{min}, x_{max}] \times [y_{min}, y_{max}]\) est initialisée comme
polygone à 4 sommets. 2. Le polygone est découpé successivement par les
4 demi-plans formant le rectangle, puis par le demi-plan interfacial :
\[\mathcal{H} = \{ (x, y) \in \mathbb{R}^2 \mid n_x x + n_y y + \alpha \le 0 \}\]
3. L'aire du polygone résultant
\(\mathcal{P} = \{ (x_k, y_k) \}_{k=1}^m\) est calculée par la formule
des lacets de Gauss (Shoelace formula) :
\[\mathcal{A}(\alpha) = \frac{1}{2} \left| \sum_{k=1}^m (x_k y_{k+1} - x_{k+1} y_k) \right| \quad \text{avec } (x_{m+1}, y_{m+1}) = (x_1, y_1)\]
4. La fonction \(\alpha \mapsto \mathcal{A}(\alpha)\) étant strictement
monotone et continue sur \([-D, D]\), la solution exacte de :
\[\mathcal{A}(\alpha) = C_{i, j} \Delta x \Delta y\] est obtenue par
recherche dichotomique (Bisection) en 40 itérations avec une tolérance
relative inférieure à \(10^{-12}\).

\subsubsection{6.9. Algorithme d'Advection Géométrique CIAM (Calcul
d'Interface Affine par
Morceaux)}\label{algorithme-dadvection-guxe9omuxe9trique-ciam-calcul-dinterface-affine-par-morceaux}

La méthode CIAM (Lafaurie et al.~1994, Scardovelli \& Zaleski 1999)
résout l'advection du fluide par découpage directionnel d'opérateurs
(Strang splitting) :
\[C^{n+1} = \mathcal{L}_y \left( \mathcal{L}_x (C^n) \right) \quad \text{en alternance avec} \quad C^{n+2} = \mathcal{L}_x \left( \mathcal{L}_y (C^{n+1}) \right)\]

Pour chaque sous-étape unidimensionnelle (par exemple suivant \(x\)) :
1. L'interface affine est reconstruite dans chaque maille donneuse
\(i\). 2. Le volume balayé à travers la face \(i+1/2\) pendant le pas
\(\Delta t\) correspond au rectangle :
\[\delta \Omega_{i+1/2} = \begin{cases} [x_{i+1/2} - u_{i+1/2}\Delta t, x_{i+1/2}] \times [y_{j-1/2}, y_{j+1/2}] & \text{si } u_{i+1/2} > 0 \\ [x_{i+1/2}, x_{i+1/2} - u_{i+1/2}\Delta t] \times [y_{j-1/2}, y_{j+1/2}] & \text{si } u_{i+1/2} < 0 \end{cases}\]
3. Le flux géométrique interfacial \(F_{x, i+1/2}\) est l'aire exacte
découpée dans ce rectangle par le segment de droite de la maille
donneuse via Sutherland-Hodgman. 4. La fraction volumique intermédiaire
est mise à jour avec correction de compressibilité 1D :
\[C_i^* = \frac{C_i^n - \frac{1}{\Delta x} (F_{x, i+1/2} - F_{x, i-1/2})}{1 - \frac{\Delta t}{\Delta x} (u_{i+1/2} - u_{i-1/2})}\]

\subsubsection{6.10. Algorithme Non Découpé Multidimensionnel de
Weymouth et Yue
(2010)}\label{algorithme-non-duxe9coupuxe9-multidimensionnel-de-weymouth-et-yue-2010}

L'algorithme de Weymouth et Yue (2010) élimine tout découpage
directionnel d'opérateur et met à jour la fraction volumique de manière
multidimensionnelle non découplée :
\[C_i^{n+1} = C_i^n - \frac{1}{\Delta x \Delta y} \left( F_{x, i+1/2} - F_{x, i-1/2} + F_{y, j+1/2} - F_{y, j-1/2} \right)\]

Pour éliminer l'erreur de double comptage dans les coins en écoulement
diagonal (\(u \neq 0, v \neq 0\)), l'algorithme calcule le flux de coin
géométrique \(V_{xy}\) sur la boîte
\([0, |u|\Delta t] \times [0, |v|\Delta t]\) et effectue une correction
symétrique :
\[F_x = F_{x0} - \frac{1}{2} V_{xy} \operatorname{sgn}(u), \quad F_y = F_{y0} + \frac{1}{2} V_{xy} \operatorname{sgn}(v)\]
où \(F_{x0}, F_{y0}\) sont les flux géométriques unidimensionnels bruts.
Enfin, une redistribution conservative locale garantit strictement la
bornitude \(C \in [0, 1]\) même sous fort cisaillement ou déformation
tourbillonnaire.

\begin{center}\rule{0.5\linewidth}{0.5pt}\end{center}

\subsection{7. Algorithmes de Résolution Algébrique et Solveurs de
Krylov}\label{algorithmes-de-ruxe9solution-alguxe9brique-et-solveurs-de-krylov}

\subsubsection{7.1. Structure Matricielle du Système Linéaire
Couplé}\label{structure-matricielle-du-systuxe8me-linuxe9aire-coupluxe9}

À chaque pas de temps, l'assemblage de l'opérateur de Navier-Stokes avec
Lagrangien Augmenté produit un système linéaire creux :
\[\mathbf{A} \mathbf{x} = \mathbf{b}\] où le vecteur d'inconnues
concatène les vitesses horizontales et verticales :
\[\mathbf{x} = \begin{pmatrix} \mathbf{u} \\ \mathbf{v} \end{pmatrix}\]

La matrice \(\mathbf{A}\) possède une structure par blocs \(2 \times 2\)
:
\[\begin{pmatrix} \mathbf{A}_{uu} & \mathbf{A}_{uv} \\ \mathbf{A}_{vu} & \mathbf{A}_{vv} \end{pmatrix} \begin{pmatrix} \mathbf{u} \\ \mathbf{v} \end{pmatrix} = \begin{pmatrix} \mathbf{b}_u \\ \mathbf{b}_v \end{pmatrix}\]

Les blocs diagonaux \(\mathbf{A}_{uu}\) et \(\mathbf{A}_{vv}\)
représentent la diffusion-convection et la pénalisation normale de la
divergence :
\[\mathbf{A}_{uu} \sim \frac{\rho}{\Delta t}\mathbf{I} - \mu \nabla^2 - r \frac{\partial^2}{\partial x^2}\]
Les blocs extradiagonaux \(\mathbf{A}_{uv}\) et \(\mathbf{A}_{vu}\)
proviennent du terme croisé de cisaillement et de la pénalisation
croisée de la divergence :
\[\mathbf{A}_{uv} \sim -(\mu + r) \frac{\partial^2}{\partial x \partial y}\]

\subsubsection{7.2. Algorithme BiCGSTAB(2) (Stabilized Bi-Conjugate
Gradient Order
2)}\label{algorithme-bicgstab2-stabilized-bi-conjugate-gradient-order-2}

Le solveur principal de COVOF 2D est le \textbf{BiCGSTAB(2)} de Sleijpen
et van der Vorst (1993). Contrairement au BiCGSTAB standard qui minimise
les résidus avec un polynôme linéaire de degré 1 à chaque pas (ce qui
provoque des stagnations lorsque le spectre comporte des valeurs propres
imaginaires conjuguées dominées par la convection), BiCGSTAB(2) utilise
un polynôme quadratique de minimisation sur deux sous-pas.

\paragraph{Algorithme Complet :}\label{algorithme-complet}

\begin{enumerate}
\def\labelenumi{\arabic{enumi}.}
\tightlist
\item
  \textbf{Initialisation} :
  \[\mathbf{r}_0 = \mathbf{P}^{-1} (\mathbf{b} - \mathbf{A} \mathbf{x}_0), \quad \tilde{\mathbf{r}} = \mathbf{r}_0, \quad \mathbf{u}_0 = \mathbf{0}, \quad \rho_0 = 1, \quad \alpha = 0, \quad \omega = 1\]
\item
  \textbf{Boucle Itérative (\(k = 0, 1, 2, \dots\))} :

  \begin{itemize}
  \tightlist
  \item
    \textbf{Étape BiCG 1} :
    \[\rho_0 \leftarrow -\omega \rho_0, \quad \rho_1 = (\tilde{\mathbf{r}}, \mathbf{r}_0), \quad \beta = \alpha \frac{\rho_1}{\rho_0}, \quad \rho_0 = \rho_1\]
    \[\mathbf{u}_0 \leftarrow \mathbf{r}_0 - \beta \mathbf{u}_0, \quad \mathbf{u}_1 = \mathbf{P}^{-1} \mathbf{A} \mathbf{u}_0\]
    \[\gamma = (\tilde{\mathbf{r}}, \mathbf{u}_1), \quad \alpha = \frac{\rho_0}{\gamma}\]
    \[\mathbf{r}_0 \leftarrow \mathbf{r}_0 - \alpha \mathbf{u}_1, \quad \mathbf{x} \leftarrow \mathbf{x} + \alpha \mathbf{u}_0, \quad \mathbf{r}_1 = \mathbf{P}^{-1} \mathbf{A} \mathbf{r}_0\]
  \item
    \textbf{Étape BiCG 2} :
    \[\rho_1 = (\tilde{\mathbf{r}}, \mathbf{r}_1), \quad \beta = \alpha \frac{\rho_1}{\rho_0}, \quad \rho_0 = \rho_1\]
    \[\mathbf{u}_0 \leftarrow \mathbf{r}_0 - \beta \mathbf{u}_0, \quad \mathbf{u}_1 \leftarrow \mathbf{r}_1 - \beta \mathbf{u}_1, \quad \mathbf{u}_2 = \mathbf{P}^{-1} \mathbf{A} \mathbf{u}_1\]
    \[\gamma = (\tilde{\mathbf{r}}, \mathbf{u}_2), \quad \alpha = \frac{\rho_0}{\gamma}\]
    \[\mathbf{x} \leftarrow \mathbf{x} + \alpha \mathbf{u}_0, \quad \mathbf{r}_0 \leftarrow \mathbf{r}_0 - \alpha \mathbf{u}_1, \quad \mathbf{r}_1 \leftarrow \mathbf{r}_1 - \alpha \mathbf{u}_2, \quad \mathbf{r}_2 = \mathbf{P}^{-1} \mathbf{A} \mathbf{r}_1\]
  \item
    \textbf{Minimisation Quadratique (MR-part)} : Résolution du système
    \(2 \times 2\) :
    \[\begin{pmatrix} (\mathbf{r}_1, \mathbf{r}_1) & (\mathbf{r}_2, \mathbf{r}_1) \\ (\mathbf{r}_2, \mathbf{r}_1) & (\mathbf{r}_2, \mathbf{r}_2) \end{pmatrix} \begin{pmatrix} \gamma_1 \\ \gamma_2 \end{pmatrix} = \begin{pmatrix} (\mathbf{r}_0, \mathbf{r}_1) \\ (\mathbf{r}_0, \mathbf{r}_2) \end{pmatrix}\]
    Mise à jour :
    \[\mathbf{x} \leftarrow \mathbf{x} + \gamma_1 \mathbf{r}_0 + \gamma_2 \mathbf{r}_1\]
    \[\mathbf{r}_0 \leftarrow \mathbf{r}_0 - (\gamma_1 \mathbf{r}_1 + \gamma_2 \mathbf{r}_2)\]
    \[\mathbf{u}_0 \leftarrow \mathbf{u}_0 - (\gamma_1 \mathbf{u}_1 + \gamma_2 \mathbf{u}_2)\]
    \[\omega = \gamma_2\]
  \item
    \textbf{Test de convergence} : Si
    \(\|\mathbf{r}_0\|_2 / \|\mathbf{b}\|_2 < \epsilon_{tol}\), arrêt.
  \end{itemize}
\end{enumerate}

\subsubsection{7.3. Algorithme GMRES(m) (Restarted Generalized Minimal
Residual)}\label{algorithme-gmresm-restarted-generalized-minimal-residual}

Pour les cas à très forte convection non symétrique, l'algorithme GMRES
avec orthogonalisation d'Arnoldi et redémarrage toutes les \(m\)
itérations (\(m = 30\) par défaut) est disponible. Il garantit la
minimisation monotone de la norme euclidienne du résidu dans le
sous-espace de Krylov \(\mathcal{K}_m(\mathbf{A}, \mathbf{r}_0)\).

\subsubsection{7.4. Algorithme du Gradient Conjugué Préconditionné
(PCG)}\label{algorithme-du-gradient-conjuguuxe9-pruxe9conditionnuxe9-pcg}

Pour l'équation de Poisson du Laplacien cartésien (Cas 3.1) où la
matrice discrète \(\mathbf{A} = -\nabla^2\) est strictement symétrique
définie positive (SPD), COVOF 2D emploie le Gradient Conjugué standard
préconditionné par Jacobi.

\subsubsection{7.5. Techniques de
Préconditionnement}\label{techniques-de-pruxe9conditionnement}

Le préconditionneur \(\mathbf{P}\) a pour rôle de condenser le spectre
des valeurs propres de \(\mathbf{P}^{-1}\mathbf{A}\) autour de
\((1, 0)\) dans le plan complexe. 1. \textbf{Jacobi (Diagonal)} :
\[\mathbf{P}_{diag} = \text{diag}(\mathbf{A}) \quad \Longrightarrow \quad P_{ii}^{-1} = \frac{1}{A_{ii}}\]
Avantage : Parallélisme parfait sans communication inter-processeurs
MPI. 2. \textbf{SSOR (Symmetric Successive Over-Relaxation)} :
\[\mathbf{P}_{SSOR} = \frac{1}{\omega(2-\omega)} (\mathbf{D} + \omega \mathbf{L}) \mathbf{D}^{-1} (\mathbf{D} + \omega \mathbf{U})\]
3. \textbf{Bloc-Jacobi MPI} : Chaque sous-domaine MPI applique une
factorisation locale ILU(0) sur son bloc interne, garantissant une
montée en charge optimale sur supercalculateurs.

\begin{center}\rule{0.5\linewidth}{0.5pt}\end{center}

\subsection{8. Stabilité Numérique et Critères
CFL}\label{stabilituxe9-numuxe9rique-et-crituxe8res-cfl}

Pour assurer la stabilité dynamique de l'intégration temporelle et la
précision physique des interfaces, le pas de temps global \(\Delta t\)
doit satisfaire simultanément trois restrictions fondamentales :

\subsubsection{8.1. Restriction Convective
CFL}\label{restriction-convective-cfl}

Elle impose que l'information hydrodynamique ne traverse pas plus d'une
fraction de maille par pas de temps :
\[\Delta t_{conv} \le CFL_{conv} \times \left( \frac{|u|_{max}}{\Delta x} + \frac{|v|_{max}}{\Delta y} \right)^{-1}\]
avec usuellement \(CFL_{conv} \le 0.5\) pour les schémas explicites
d'interface VOF.

\subsubsection{8.2. Restriction Visqueuse}\label{restriction-visqueuse}

Dans COVOF 2D, l'opérateur visqueux étant discrétisé de façon totalement
implicite dans le Lagrangien Augmenté, la condition de stabilité
classique de diffusion explicite :
\[\Delta t_{visc} \le \frac{\rho \Delta x^2 \Delta y^2}{2 \mu (\Delta x^2 + \Delta y^2)}\]
est \textbf{inconditionnellement contournée}. Le solveur reste stable
même pour des pas de temps \(\Delta t \gg \Delta t_{visc}\).

\subsubsection{8.3. Restriction Capillaire de
Brackbill}\label{restriction-capillaire-de-brackbill}

La propagation d'ondes capillaires à l'interface liquide/gaz impose la
condition de dispersion de Brackbill (1992) :
\[\Delta t_\sigma \le \sqrt{\frac{(\rho_1 + \rho_2) \min(\Delta x, \Delta y)^3}{4 \pi \sigma}}\]

Si cette condition est violée, des ondes de surface artificielles de
haute fréquence s'amplifient exponentiellement, conduisant à la
désintégration numérique de la goutte ou de la bulle.

\subsubsection{8.4. Algorithme d'Adaptation Automatique du Pas de
Temps}\label{algorithme-dadaptation-automatique-du-pas-de-temps}

À chaque pas de calcul \(n\), le nouveau pas de temps \(\Delta t^{n+1}\)
est réévalué dynamiquement selon :
\[\Delta t^{n+1} = \min \left( \Delta t_{max}, \, \Delta t_{conv}, \, \Delta t_\sigma, \, 1.1 \times \Delta t^n \right)\]

Le facteur \(1.1\) limite l'accroissement maximal par pas à \(+10\%\)
afin de prévenir les chocs numériques transitoires.

\begin{center}\rule{0.5\linewidth}{0.5pt}\end{center}

\subsection{9. Nomenclature et Symboles
Mathématiques}\label{nomenclature-et-symboles-mathuxe9matiques}

{\def\LTcaptype{none} % do not increment counter
\begin{longtable}[]{@{}
  >{\raggedright\arraybackslash}p{(\linewidth - 4\tabcolsep) * \real{0.3333}}
  >{\raggedright\arraybackslash}p{(\linewidth - 4\tabcolsep) * \real{0.3333}}
  >{\raggedright\arraybackslash}p{(\linewidth - 4\tabcolsep) * \real{0.3333}}@{}}
\toprule\noalign{}
\begin{minipage}[b]{\linewidth}\raggedright
Symbole
\end{minipage} & \begin{minipage}[b]{\linewidth}\raggedright
Grandeur Physique
\end{minipage} & \begin{minipage}[b]{\linewidth}\raggedright
Unité SI
\end{minipage} \\
\midrule\noalign{}
\endhead
\bottomrule\noalign{}
\endlastfoot
\(\mathbf{u} = (u, v)^T\) & Champ de vitesse eulérien &
\(\text{m/s}\) \\
\(p\) & Pression hydrodynamique & \(\text{Pa} = \text{N/m}^2\) \\
\(T\) & Température thermodynamique & \(\text{K}\) ou
\(^{\circ}\text{C}\) \\
\(C\) & Fraction volumique de couleur (VOF) & adimensionnel
\([0, 1]\) \\
\(\rho\) & Masse volumique & \(\text{kg/m}^3\) \\
\(\mu\) & Viscosité dynamique &
\(\text{Pa}\cdot\text{s} = \text{kg/(m}\cdot\text{s)}\) \\
\(\mathbf{g}\) & Accélération de la pesanteur & \(\text{m/s}^2\) \\
\(\sigma\) & Coefficient de tension de surface & \(\text{N/m}\) \\
\(\kappa\) & Courbure locale de l'interface & \(\text{m}^{-1}\) \\
\(\mathbf{n}\) & Vecteur normal unitaire sortant & adimensionnel \\
\(r\) & Paramètre d'augmentation du Lagrangien &
\(\text{kg/(m}^3\cdot\text{s)}\) \\
\(\beta\) & Coefficient de dilatation thermique & \(\text{K}^{-1}\) \\
\(k\) & Conductivité thermique & \(\text{W/(m}\cdot\text{K)}\) \\
\(c_p\) & Capacité thermique isobare & \(\text{J/(kg}\cdot\text{K)}\) \\
\(K\) & Perméabilité de Brinkman & \(\text{m}^2\) \\
\(Nu, \overline{Nu}\) & Nombre de Nusselt local et moyen &
adimensionnel \\
\(Re\) & Nombre de Reynolds \(\rho U L / \mu\) & adimensionnel \\
\(Ra\) & Nombre de Rayleigh \(g \beta \Delta T L^3 / (\nu \alpha)\) &
adimensionnel \\
\(Pr\) & Nombre de Prandtl \(\nu / \alpha\) & adimensionnel \\
\end{longtable}
}

\begin{center}\rule{0.5\linewidth}{0.5pt}\end{center}

\subsection{10. Références
Bibliographiques}\label{ruxe9fuxe9rences-bibliographiques}

\begin{enumerate}
\def\labelenumi{\arabic{enumi}.}
\tightlist
\item
  \textbf{Harlow, F. H., \& Welch, J. E. (1965)}. \emph{Numerical
  calculation of time-dependent viscous incompressible flow of fluid
  with free surface}. Physics of Fluids, 8(12), 2182-2189.
\item
  \textbf{Fortin, M., \& Glowinski, R. (1983)}. \emph{Augmented
  Lagrangian methods: applications to the numerical solution of
  boundary-value problems}. North-Holland Publishing Company.
\item
  \textbf{Caltagirone, J. P., \& Vincent, S. (2001)}. \emph{Fluid-fluid
  interface tracking by a projection method and direct calculation of
  surface tension}. Computers \& Fluids, 30(7-8), 913-930.
\item
  \textbf{Brackbill, J. U., Kothe, D. B., \& Zemach, C. (1992)}. \emph{A
  continuum method for modeling surface tension}. Journal of
  Computational Physics, 100(2), 335-354.
\item
  \textbf{François, M. M., Cummins, S. J., Dendy, E. D., et al.~(2006)}.
  \emph{A balanced-force algorithm for continuous and sharp surface
  tension models within a volume tracking framework}. Journal of
  Computational Physics, 213(1), 141-173.
\item
  \textbf{Sleijpen, G. L., \& van der Vorst, H. A. (1993)}.
  \emph{BiCGSTAB(2) for linear equations involving unsymmetric matrices
  with complex spectrum}. Electronic Transactions on Numerical Analysis,
  1, 11-32.
\item
  \textbf{Laboratoire TREFLE / I2M (2007)}. \emph{Cahier de Validation
  et Vérification CFD 2D}. UMR CNRS 5295, Université de Bordeaux.
\item
  \textbf{Vincent, S. (1999)}. \emph{Modélisation d'écoulements
  turbulents multiphasiques en mécanique des fluides numérique}. Thèse
  de doctorat, Université Bordeaux I.
\end{enumerate}

\begin{center}\rule{0.5\linewidth}{0.5pt}\end{center}

\subsection{11. Développements Mathématiques Avancés et Preuves de
Convergence}\label{duxe9veloppements-mathuxe9matiques-avancuxe9s-et-preuves-de-convergence}

\subsubsection{11.1. Forme Algébrique Développée des Coefficients
Stencils}\label{forme-alguxe9brique-duxe9veloppuxe9e-des-coefficients-stencils}

Soit le système linéaire discrétisé assemblé au pas de temps \(n+1\) sur
la grille décalée MAC. Pour chaque degré de liberté \(u_{i, j}\) associé
à la face verticale située entre les cellules \((i-1, j)\) et \((i, j)\)
:

\[a_{P, u} u_{i, j} + a_{E, u} u_{i+1, j} + a_{W, u} u_{i-1, j} + a_{N, u} u_{i, j+1} + a_{S, u} u_{i, j-1} + \sum_{c \in \{NE, NW, SE, SW\}} a_{c, v} v_c = b_{u, i, j}\]

Les coefficients matriciels s'explicitent analytiquement comme suit :

\paragraph{\texorpdfstring{1. Terme Diagonal Principal
\(a_{P, u}\)}{1. Terme Diagonal Principal a\_\{P, u\}}}\label{terme-diagonal-principal-a_p-u}

\[a_{P, u} = \frac{\rho_{i-1, j} + \rho_{i, j}}{2 \Delta t} + 2 \left( \frac{\mu_{i, j} + \mu_{i-1, j}}{\Delta x^2} \right) + \left( \frac{\mu_{node, i, j+1} + \mu_{node, i, j}}{\Delta y^2} \right) + \frac{2 r_{al}}{\Delta x^2} + \frac{\mu_{i, j}}{K_{solid}} \delta_{solid}\]
où \(\delta_{solid} = 1\) si la face est située à l'intérieur d'un
obstacle solide pénalisé par Brinkman, et \(0\) sinon.

\paragraph{\texorpdfstring{2. Termes de Voisinage Axiaux \(a_{E, u}\) et
\(a_{W, u}\)}{2. Termes de Voisinage Axiaux a\_\{E, u\} et a\_\{W, u\}}}\label{termes-de-voisinage-axiaux-a_e-u-et-a_w-u}

\[a_{E, u} = - \frac{2 \mu_{i, j}}{\Delta x^2} - \frac{r_{al}}{\Delta x^2}\]
\[a_{W, u} = - \frac{2 \mu_{i-1, j}}{\Delta x^2} - \frac{r_{al}}{\Delta x^2}\]

\paragraph{\texorpdfstring{3. Termes de Voisinage Transversaux
\(a_{N, u}\) et
\(a_{S, u}\)}{3. Termes de Voisinage Transversaux a\_\{N, u\} et a\_\{S, u\}}}\label{termes-de-voisinage-transversaux-a_n-u-et-a_s-u}

\[a_{N, u} = - \frac{\mu_{node, i, j+1}}{\Delta y^2}\]
\[a_{S, u} = - \frac{\mu_{node, i, j}}{\Delta y^2}\]

\paragraph{\texorpdfstring{4. Termes de Couplage Croisé Vitesse-Vitesse
\(a_{c, v}\)}{4. Termes de Couplage Croisé Vitesse-Vitesse a\_\{c, v\}}}\label{termes-de-couplage-croisuxe9-vitesse-vitesse-a_c-v}

Ces termes proviennent du cisaillement visqueux
\(\frac{\partial}{\partial y}\left( \mu \frac{\partial v}{\partial x} \right)\)
et de la pénalisation croisée de la divergence
\(-r_{al} \frac{\partial}{\partial x}\left( \frac{\partial v}{\partial y} \right)\)
:
\[a_{NE, v} = - \frac{\mu_{node, i, j+1} + r_{al}}{\Delta x \Delta y}, \quad a_{NW, v} = + \frac{\mu_{node, i, j+1} + r_{al}}{\Delta x \Delta y}\]
\[a_{SE, v} = + \frac{\mu_{node, i, j} + r_{al}}{\Delta x \Delta y}, \quad a_{SW, v} = - \frac{\mu_{node, i, j} + r_{al}}{\Delta x \Delta y}\]

\subsubsection{11.2. Preuve Mathématique de la Convergence de
l'Algorithme
d'Uzawa}\label{preuve-mathuxe9matique-de-la-convergence-de-lalgorithme-duzawa}

Considérons le problème de point-selle sous forme d'opérateurs continus
ou matriciels discrets :
\[\mathbf{A}_r \mathbf{u} + \mathbf{B}^T p = \mathbf{f}\]
\[\mathbf{B} \mathbf{u} = \mathbf{0}\] où
\(\mathbf{A}_r = \mathbf{A}_0 + r \mathbf{B}^T \mathbf{B}\) est
strictement symétrique définie positive (en régime de Stokes ou sous
forte dissipation visqueuse).

L'itération d'Uzawa pour la pression s'écrit :
\[p^{k+1} = p^k - \rho_u \mathbf{B} \mathbf{u}^{k+1}\] avec
\(\mathbf{u}^{k+1} = \mathbf{A}_r^{-1} (\mathbf{f} - \mathbf{B}^T p^k)\).
En soustrayant la pression exacte \(p^*\), l'erreur
\(e_p^k = p^k - p^*\) vérifie la relation de récurrence linéaire :
\[e_p^{k+1} = \left( \mathbf{I} - \rho_u \mathbf{B} \mathbf{A}_r^{-1} \mathbf{B}^T \right) e_p^k\]

\paragraph{Théorème de Convergence}\label{thuxe9oruxe8me-de-convergence}

Si le pas d'Uzawa est choisi égal au paramètre d'augmentation
(\(\rho_u = r\)), la matrice d'itération
\(\mathbf{M}_r = \mathbf{I} - r \mathbf{B} \mathbf{A}_r^{-1} \mathbf{B}^T\)
vérifie :
\[\|\mathbf{M}_r\|_2 = \frac{1}{1 + r \lambda_{min}(\mathbf{B} \mathbf{A}_0^{-1} \mathbf{B}^T)} < 1\]
où \(\lambda_{min}\) désigne la plus petite valeur propre non nulle du
complément de Schur de Stokes. Par conséquent, lorsque \(r \to \infty\),
le rayon spectral satisfait asymptotiquement :
\[\rho(\mathbf{M}_r) \approx \frac{C}{r} \longrightarrow 0\] Ce résultat
fondamental démontre que \textbf{la convergence d'Uzawa devient
instantanée (en 1 à 2 itérations)} lorsque le paramètre d'augmentation
\(r\) est suffisamment grand, validant théoriquement l'architecture sans
inversion de Poisson de COVOF 2D.

\subsubsection{11.3. Estimation d'Erreur de la Pénalisation de
Brinkman}\label{estimation-derreur-de-la-puxe9nalisation-de-brinkman}

Soit \(\mathbf{u}\) la solution exacte des équations de Navier-Stokes
dans le domaine fluide \(\Omega_f\) avec condition d'adhérence
\(\mathbf{u} = \mathbf{0}\) sur l'interface solide
\(\Gamma_s = \partial \Omega_s\). Soit \(\mathbf{u}_K\) la solution de
l'équation pénalisée dans le domaine étendu
\(\Omega = \Omega_f \cup \Omega_s\) :
\[\rho \left( \frac{\partial \mathbf{u}_K}{\partial t} + (\mathbf{u}_K \cdot \nabla)\mathbf{u}_K \right) - \nabla \cdot (2\mu \mathbf{D}(\mathbf{u}_K)) + \nabla p_K + \frac{\mu}{K(\mathbf{x})} \mathbf{u}_K = \mathbf{f}\]
avec \(K(\mathbf{x}) = +\infty\) dans \(\Omega_f\) et
\(K(\mathbf{x}) = K_{solid}\) dans \(\Omega_s\).

\paragraph{Théorème (Angot, Bruneau \& Fabrie
1999)}\label{thuxe9oruxe8me-angot-bruneau-fabrie-1999}

Pour un écoulement incompressible régulier, la solution pénalisée
\(\mathbf{u}_K\) converge fortement vers la solution exacte
\(\mathbf{u}\) dans \(L^2(0, T; \mathbf{H}^1(\Omega_f))\) selon
l'estimation :
\[\|\mathbf{u} - \mathbf{u}_K\|_{L^2(0, T; \mathbf{L}^2(\Omega_f))} \le C_1 \sqrt{K_{solid}}\]
\[\|\mathbf{u}_K\|_{L^2(0, T; \mathbf{L}^2(\Omega_s))} \le C_2 K_{solid}\]

En choisissant \(K_{solid} = 10^{-40}\text{ m}^2\) dans COVOF 2D,
l'erreur de modélisation géométrique de l'obstacle :
\[\|\mathbf{u} - \mathbf{u}_K\| \sim 10^{-20}\] est
\textbf{incommensurablement plus faible que l'erreur de discrétisation
spatiale} \(O(\Delta x^2) \sim 10^{-4}\), garantissant une précision
géométrique optimale sans maillage déformé.

\subsubsection{11.4. Diagramme de Sweby et Préservation de la Monotonie
(TVD)}\label{diagramme-de-sweby-et-pruxe9servation-de-la-monotonie-tvd}

Pour le transport d'une fraction volumique \(C \in [0, 1]\) gouverné par
\(\frac{\partial C}{\partial t} + u \frac{\partial C}{\partial x} = 0\),
la condition de variation totale non croissante (TVD) de Harten (1983)
stipule :
\[TV(C^{n+1}) \le TV(C^n) \quad \text{avec} \quad TV(C) = \sum_i |C_{i+1} - C_i|\]

Dans le plan \((\theta, \Phi(\theta))\) de Sweby, la région de stabilité
TVD est délimitée par :
\[\Phi(\theta) \le \min(2\theta, 2) \quad \text{pour } \theta > 0, \quad \Phi(\theta) = 0 \quad \text{pour } \theta \le 0\]

\begin{verbatim}
   Phi(theta)
       ▲
     2 ┼            ┌──────────────────────── Limiteur SUPERBEE (Roe)
       │           /│
       │          / │
     1 ┼─────────/──┼──────────────────────── Schéma d'Ordre 2 Centré
       │        /   │
       │       /    │                        Région TVD de Sweby
       │      /     │
     0 ┼─────┴──────┴────────────────────────► theta
       0     0.5    1     2
\end{verbatim}

Le limiteur \textbf{Superbee de Roe} utilisé dans COVOF 2D longe la
frontière supérieure de la région de Sweby :
\[\Phi_{Superbee}(\theta) = \max \left( 0, \max(\min(2\theta, 1), \min(\theta, 2)) \right)\]
Cette propriété confère au solveur de transport d'interface sa capacité
démontrée à maintenir les sauts de fraction volumique sur 1 à 2 mailles
sans diffusion ni dispersion oscillatoire.

\end{document}
